% ============================================================
%  Week 8 — Mediation and Front-Door Identification
%  Chapter 8 in the master book
%  Compile standalone:  pdflatex chapter8_new.tex
%  Compile via master:  pdflatex causal_inference_master.tex
% ============================================================
\documentclass[master.tex]{subfiles}

\begin{document}

% ------------------------------------------------------------
%  Title block (standalone) / chapter heading (master)
% ------------------------------------------------------------
\ifSubfilesClassLoaded{%
  \begin{center}
    {\LARGE\bfseries\color{isublue}
     Chapter 8: Mediation and Front-Door Identification}\\[6pt]
    {\large Theory and Methods in Causal Inference}\\[4pt]
    {\normalsize Department of Statistics, Iowa State University}
  \end{center}
  \bigskip
  \hrule height 1.2pt \relax
  \smallskip
}{%
  \chapter{Mediation and Front-Door Identification}
}

% ------------------------------------------------------------
%  Learning objectives
% ------------------------------------------------------------
\begin{tcolorbox}[defstyle, title={Learning Objectives}]
By the end of this chapter, students should be able to:
\begin{enumerate}
  \item Explain the distinction between the total causal effect of $T$ on
    $Y$ and a pathway-specific effect that operates through a mediator $M$,
    and describe why this distinction matters scientifically.
  \item Draw the prototype mediation DAG, write its structural equations,
    and identify the direct and indirect pathways.
  \item Define the controlled direct effect (CDE) using the do-operator,
    identify it via the back-door formula for the joint intervention
    $(T, M)$, and explain why it depends on the fixed level $m$.
  \item Define the natural direct and indirect effects (NDE, NIE) using
    the potential outcomes notation $Y(t, M(t'))$, state the
    NDE$+$NIE$=$TE decomposition, and explain why these quantities involve
    cross-world counterfactuals.
  \item State the four sequential ignorability assumptions for
    identification of natural effects, write the mediation formula, and
    identify which assumption is violated when there is an unmeasured
    treatment-induced mediator--outcome confounder.
  \item Set up the Baron--Kenny three-equation system, derive the product
    and difference formulas for the indirect effect, and explain why the
    decomposition fails in nonlinear or interaction models.
  \item State the three front-door conditions, derive the front-door
    formula using the do-calculus, and explain why the front-door graph
    enables identification despite unobserved $T$--$Y$ confounding.
  \item Contrast mediation analysis and instrumental variables on the
    dimensions of variable position, identification goal, and key
    assumption.
\end{enumerate}
\end{tcolorbox}

% ------------------------------------------------------------
%  Local styles
% ------------------------------------------------------------
% (\E defined in master.tex)
% (remarkstyle and remark environment are defined in causal_inference_master.tex)

% ============================================================
\section{Motivation: Mechanisms}
\label{w8:sec:motivation}
% ============================================================

The identification results of Chapters~5--7 all answer the same
question: \emph{what is the total causal effect of $T$ on $Y$?}
Mediation analysis asks a finer question: \emph{through what mechanism
does that effect operate?}

More concretely, the total effect of $T$ on $Y$ may flow along multiple
causal pathways.  Some of this effect passes through an intermediate
variable $M$ --- the \emph{mediator} --- along the path $T \to M \to Y$.
The remainder flows directly along $T \to Y$, bypassing the mediator
entirely.  Decomposing the total effect into these two components is the
goal of mediation analysis.

This decomposition matters for scientific and policy reasons.  In a
clinical trial of a behavioral intervention ($T$) on depression ($Y$),
a researcher may want to know how much of the benefit operates through
improved sleep quality ($M$) versus other pathways --- because if sleep
is the main channel, targeting sleep directly may be a more efficient
intervention.  In an economics study of education ($T$) on wages ($Y$),
how much operates through occupation ($M$) versus cognitive skills?
The answer determines whether a policy should target educational
attainment or occupational access.

The challenge is that mediators are \emph{post-treatment variables}:
they are affected by the treatment, and may themselves be confounded
with the outcome.  Conditioning on a post-treatment variable creates
exactly the collider and selection-bias problems studied in Chapters~2
and~3.  A naive approach --- simply including $M$ as a covariate in a
regression of $Y$ on $T$ --- conflates adjustment with mediation and
can introduce bias even in a randomized experiment.

\medskip\noindent This chapter also develops the \emph{front-door
criterion}, a distinct identification strategy that uses the mediation
structure of the DAG to identify causal effects even when treatment and
outcome are confounded by an unobserved variable.  This makes
mediation analysis relevant not only to mechanism research but also to
the core identification problem of earlier chapters.

% ============================================================
\section{The Mediation DAG}
\label{w8:sec:dag}
% ============================================================

\subsection{The Prototype Graph}

Throughout this chapter we work with the following prototype graph,
shown in Figure~\ref{w8:fig:mediation-dag}.

\begin{figure}[h]
\centering
\begin{tikzpicture}[node distance=1.9cm]
  \node[node]  (T) at (0,0)      {$T$};
  \node[node]  (M) at (2.4,0)    {$M$};
  \node[node]  (Y) at (4.8,0)    {$Y$};
  \node[node]  (X) at (2.4, 1.5) {$\mathbf{X}$};
  \node[unode] (U) at (2.4,-1.5) {$U$};

  \draw[edge]  (T) -- (M);
  \draw[edge]  (M) -- (Y);
  \draw[edge]  (T) to[bend right=20] (Y);
  \draw[edge]  (X) -- (T);
  \draw[edge]  (X) -- (M);
  \draw[edge]  (X) -- (Y);
  \draw[dedge] (U) -- (T);
  \draw[dedge] (U) -- (Y);
  \node[font=\small\itshape\color{darkgrey}, below=0.12cm of T] {treatment};
  \node[font=\small\itshape\color{darkgrey}, below=0.12cm of M] {mediator};
  \node[font=\small\itshape\color{darkgrey}, below=0.12cm of Y] {outcome};
  \node[font=\small\itshape\color{darkgrey}, above=0.10cm of X] {covariates};
  \node[font=\small\itshape\color{darkgrey}, below=0.12cm of U] {unobserved};
\end{tikzpicture}
\caption{The prototype mediation DAG.  The causal effect of $T$ on $Y$
operates through two pathways: the direct path $T \to Y$ and the
indirect path $T \to M \to Y$ via the mediator.  The unobserved
variable $U$ confounds the treatment--outcome relationship;
$\mathbf{X}$ denotes observed pre-treatment covariates.}
\label{w8:fig:mediation-dag}
\end{figure}

\noindent The graph encodes two causal pathways:
\begin{itemize}
  \item \textbf{Direct pathway}: $T \to Y$.  The treatment affects the
    outcome without passing through the mediator.
  \item \textbf{Indirect pathway}: $T \to M \to Y$.  The treatment first
    shifts the mediator, which in turn shifts the outcome.
\end{itemize}

\subsection{Structural Equations}

The prototype graph corresponds to the following structural equation
model (SEM):
\begin{align}
  M &= f_M(T,\, \mathbf{X},\, U_M), \label{w8:eq:sem-m} \\
  Y &= f_Y(T,\, M,\, \mathbf{X},\, U_Y), \label{w8:eq:sem-y}
\end{align}
where $U_M$ and $U_Y$ are structural errors that may be correlated with
each other and with the unobserved confounder $U$.  The functions
$f_M$ and $f_Y$ are unrestricted at this stage; the linear special case
is treated in §\ref{w8:sec:linear}.

\subsection{What Makes Mediation Harder Than Total Effect Estimation}

The key difficulty is that $M$ is a \emph{post-treatment} variable.
This creates two interrelated problems.

\medskip\noindent\textbf{Collider bias.}  Conditioning on $M$ can open
collider paths.  Suppose $U \to T$ and $V \to M$ and $V \to Y$, with
$V$ unobserved.  The path $T \to M \leftarrow V \to Y$ is blocked when
$M$ is not conditioned on, but opens as soon as $M$ is included as a
covariate, inducing spurious $T$--$Y$ association through $V$.  This
is precisely why naively regressing $Y$ on $(T, M)$ does not isolate
the direct effect.

\medskip\noindent\textbf{Mediator--outcome confounding.}  Even when
treatment is randomized --- eliminating confounding on the $T$ side ---
the mediator $M$ is never randomized.  An unobserved variable $V$ with
$V \to M$ and $V \to Y$ creates a back-door path from $M$ to $Y$ that
randomization of $T$ does not close.  This is the central challenge of
mediation analysis.

\medskip\noindent These difficulties motivate the rest of the chapter.
Sections~\ref{w8:sec:cde}--\ref{w8:sec:id-natural} address the
decomposition problem under suitable assumptions on confounding.
Section~\ref{w8:sec:frontdoor} shows that, in a special case where the
direct $T \to Y$ edge is absent and $U$ does not affect $M$, the
mediation structure can be \emph{exploited} to identify the total causal
effect even when $U$ is wholly unobserved and no back-door adjustment
set exists.

\begin{warn}{Post-Treatment Variables Are Dangerous to Condition On}
\label{w8:warn:post-treatment}
The mediator $M$ is caused by the treatment $T$.  This single fact
makes every operation involving $M$ structurally different from
operations involving pre-treatment covariates $\mathbf{X}$.

\medskip\noindent Conditioning on a post-treatment variable in a
regression or matching procedure is \emph{not} a neutral act.  It can
open collider paths that were previously blocked, introduce
selection bias, and produce estimates of neither the total effect nor
any well-defined direct effect.  The appropriate response is not to
avoid $M$ entirely --- mediation analysis requires reasoning about $M$
--- but to be precise about \emph{which estimand} one is targeting and
\emph{which structural assumptions} justify the operation being
performed.

\medskip\noindent This chapter studies three regimes in which
conditioning on or marginalizing over $M$ is justified:
\begin{enumerate}
  \item \textbf{CDE} (§\ref{w8:sec:cde}): both $T$ and $M$ are
    intervened on simultaneously via $\doop(T, M)$, so no conditioning
    on an observed $M$ occurs.
  \item \textbf{NDE/NIE} (§\ref{w8:sec:nie}--\ref{w8:sec:id-natural}):
    $M$ is marginalized over using a distribution from a different
    treatment arm, justified by sequential ignorability.
  \item \textbf{Front-door} (§\ref{w8:sec:frontdoor}): $M$ is summed
    over in the front-door formula, justified by the graphical
    conditions of the front-door criterion.
\end{enumerate}
Outside these three regimes, conditioning on a post-treatment variable
should be treated as an error until proven otherwise.
\end{warn}

% ============================================================
\section{Total Causal Effect}
\label{w8:sec:te}
% ============================================================

Before decomposing the causal effect into direct and indirect
components, we fix the reference quantity: the \emph{total causal
effect} of $T$ on $Y$.

\begin{defn}{Total Effect}
\label{w8:defn:te}
The \textbf{total effect} (TE) of $T$ on $Y$ is:
\[
  \mathrm{TE}
  = \E\!\bigl[Y(1) - Y(0)\bigr]
  = \E\!\bigl[Y \mid \doop(T{=}1)\bigr]
  - \E\!\bigl[Y \mid \doop(T{=}0)\bigr].
\]
\end{defn}

\noindent The total effect captures the combined impact of all causal
pathways from $T$ to $Y$ --- both the direct path $T \to Y$ and the
indirect path $T \to M \to Y$.  It is identified by the back-door
formula whenever $\mathbf{X}$ satisfies the back-door criterion for the
effect of $T$ on $Y$:
\[
  \mathrm{TE}
  = \sum_{\mathbf{x}}
    \bigl[\E(Y \mid T{=}1, \mathbf{X}{=}\mathbf{x})
        - \E(Y \mid T{=}0, \mathbf{X}{=}\mathbf{x})\bigr]
    P(\mathbf{X}{=}\mathbf{x}).
\]
The goal of mediation analysis is to decompose this total into a
\emph{direct} component and an \emph{indirect} component.  Sections
\ref{w8:sec:cde}--\ref{w8:sec:nie} introduce two decompositions, each
formalising a different notion of what ``direct'' and ``indirect''
mean.

% ============================================================
\section{Controlled Direct Effect}
\label{w8:sec:cde}
% ============================================================

\subsection{Definition}

The do-calculus provides the cleanest definition of the direct effect:
intervene on both $T$ \emph{and} $M$ simultaneously, fixing $M$ at a
specified level $m$.  Any remaining effect of $T$ on $Y$ must then flow
exclusively through the direct edge $T \to Y$.

\begin{defn}{Controlled Direct Effect (CDE)}
\label{w8:defn:cde}
The \textbf{controlled direct effect} of changing $T$ from $t'$ to $t$
while fixing $M = m$ by intervention is:
\begin{equation}
  \mathrm{CDE}(m)
  \;=\; \E\!\left[Y \mid \doop(T{=}1),\, \doop(M{=}m)\right]
       - \E\!\left[Y \mid \doop(T{=}0),\, \doop(M{=}m)\right].
  \label{w8:eq:cde}
\end{equation}
\end{defn}

Because~\eqref{w8:eq:cde} involves two simultaneous interventions, it
corresponds to the mutilated graph $\Gcal_{\overline{TM}}$ in which all
edges into both $T$ and $M$ are deleted.  Fixing $M = m$ for everyone
shuts down the indirect pathway $T \to M \to Y$: any remaining
$T$-to-$Y$ effect flows only through the direct edge.

\begin{remark}
The CDE depends on the level $m$ at which $M$ is fixed.  In general,
the direct effect may vary across values of $m$ --- this is
\emph{effect modification by the mediator}.  In a linear additive
model, $\mathrm{CDE}(m) = \tau'$ for all $m$, which is a special
property of linearity.  When $T$ and $M$ interact in their effect on
$Y$, the CDEs at different values of $m$ differ, and no single number
summarizes the direct effect.
\end{remark}

\subsection{Identification of the CDE}

Since the CDE involves a joint intervention $\doop(T, M)$, its
identification reduces to applying the back-door criterion to the pair
$(T, M)$ jointly.

\begin{thm}{Identification of the CDE}
\label{w8:thm:cde}
Suppose $\mathbf{X}$ satisfies the back-door criterion for the joint
intervention $\doop(T, M)$ on $Y$: (i)~$\mathbf{X}$ contains no
descendant of $T$ or $M$, and (ii)~$\mathbf{X}$ blocks all back-door
paths into $(T, M)$.  Then:
\begin{equation}
  \E[Y \mid \doop(T{=}t),\, \doop(M{=}m)]
  = \sum_{\mathbf{x}} \E[Y \mid t,\, m,\, \mathbf{x}]\, P(\mathbf{x}).
  \label{w8:eq:cde-id}
\end{equation}
\end{thm}

The identification formula~\eqref{w8:eq:cde-id} is the standard
back-door adjustment formula applied to the pair $(T, M)$ as a joint
treatment.  No assumptions beyond the back-door criterion are needed.
In Figure~\ref{w8:fig:mediation-dag}, the mutilated graph
$\Gcal_{\overline{TM}}$ deletes all edges into both $T$ and $M$,
removing the confounding paths through $U$; $\mathbf{X}$ then blocks
the remaining back-door path $T \leftarrow \mathbf{X} \to Y$.

\subsection{Physical Manipulability and the Meaning of the CDE}

The CDE answers the question: what is the effect of $T$ on $Y$ when
the mediator is \emph{prevented} from changing?  This question has a
clean causal answer precisely because $\doop(M{=}m)$ is a genuine
intervention: it severs all edges into $M$, placing $M$ at $m$
regardless of what $T$ does.

\begin{remark}
The CDE is only scientifically meaningful when an intervention to fix
$M$ at a specified level $m$ is \emph{physically realizable} in the
study context.  Examples where this is plausible include a drug trial
where both the assigned treatment and a downstream biomarker can be
externally controlled, or an experiment where mediator levels are
directly assigned.  In many substantive settings, however, the mediator
cannot be independently manipulated:
\begin{itemize}
  \item In an education study, one cannot intervene to hold occupation
    fixed while varying years of schooling.
  \item In a behavioral trial, one cannot set sleep quality to a
    predetermined level while varying the intervention.
\end{itemize}
When the mediator is not cleanly manipulable, the CDE remains
mathematically defined but its policy interpretation is strained.  In
those cases the natural direct effect (§\ref{w8:sec:nie}), which asks
what would happen if $M$ were allowed to take its natural value under
the reference treatment, may be more interpretively coherent ---
though at the cost of stronger identification assumptions.
\end{remark}

\subsection{The CDE Does Not Have a Complementary Indirect Effect}

A common misconception is that the CDE and some complementary
``controlled indirect effect'' sum to the total effect, analogous to
the NDE$+$NIE$=$TE decomposition.  This is not generally true.

Fixing $M = m$ by intervention removes the pathway $T \to M \to Y$
from the mutilated graph, but the residual quantity
$\mathrm{TE} - \mathrm{CDE}(m)$ is not a causal effect in its own
right.  It depends on $m$, it has no clean do-calculus expression
corresponding to ``the indirect pathway,'' and it cannot be
interpreted as the effect that flows through $M$ in any natural
sense.  Formally:
\begin{equation}
  \mathrm{TE} \;-\; \mathrm{CDE}(m)
  \;=\;
  \E\!\left[Y \mid \doop(T{=}1)\right]
  - \E\!\left[Y \mid \doop(T{=}1),\,\doop(M{=}m)\right]
  - \left(\E\!\left[Y \mid \doop(T{=}0)\right]
  - \E\!\left[Y \mid \doop(T{=}0),\,\doop(M{=}m)\right]\right),
  \label{w8:eq:te-minus-cde}
\end{equation}
which mixes the effect of intervening on $M$ under $T=1$ and $T=0$ in
a way that has no simple pathway interpretation.  The decomposition
TE $=$ NDE $+$ NIE holds because the natural effects involve
cross-world counterfactuals that are constructed precisely to be
complementary; the CDE has no analogous cross-world partner.

\medskip\noindent The practical implication is that if the research
goal is to decompose the total effect into direct and indirect
components, the correct estimands are the NDE and NIE, not
$\mathrm{CDE}(m)$ and $\mathrm{TE} - \mathrm{CDE}(m)$.  The CDE is
the right estimand when the question is specifically about the effect
of $T$ with $M$ controlled at a given level --- a different and more
limited question.

Natural effects, defined in the next section, capture the more natural
decomposition notion.

\begin{warn}{Natural Effects Are a Strictly Harder Estimand Than the CDE}
\label{w8:warn:crossworld}
The step from the CDE to the NDE and NIE is not merely a notational
upgrade --- it is a conceptual escalation with serious identification
consequences.

\medskip\noindent\textbf{The CDE is a single-world estimand.}  The
quantity $\E[Y \mid \doop(T{=}t), \doop(M{=}m)]$ involves a single
joint intervention.  Every individual lives in one hypothetical world
where both $T$ and $M$ are externally fixed.  The do-operator handles
this directly, and identification reduces to the standard back-door
criterion for the pair $(T, M)$.

\medskip\noindent\textbf{The NDE and NIE are cross-world estimands.}
The nested counterfactual $Y(t, M(t'))$ requires an individual to
simultaneously inhabit two worlds: the world where $T = t$ (which
determines $Y$) and the world where $T = t'$ (which determines what
$M$ would \emph{naturally} be).  When $t \neq t'$, these are logically
distinct states of the world.  No single experimental intervention can
realize both at once for the same unit.  The do-operator alone cannot
express this quantity.

\medskip\noindent\textbf{The identification price is steep.}  The CDE
requires only the back-door criterion --- an assumption about observed
confounding that can in principle be satisfied by design.  The NDE and
NIE require four sequential ignorability assumptions, including
Assumption~4 (no treatment-induced mediator--outcome confounder), which
\emph{cannot} be satisfied by randomizing $T$ and cannot be verified
from data.  This is why natural effects remain controversial in
settings where the treatment-induced confounder assumption is not
credible on subject-matter grounds.

\medskip\noindent\textbf{Practical guidance.}  When the goal is a
pathway-specific effect and the no-treatment-induced-confounder
assumption is plausible, NDE and NIE are the right estimands.  When
that assumption is suspect, the CDE --- or a sensitivity analysis
around it --- is the safer choice.
\end{warn}

% ============================================================
\section{Natural Direct and Indirect Effects}
\label{w8:sec:nie}
% ============================================================

\subsection{Motivation: Cross-World Counterfactuals}

The CDE fixes the mediator by external intervention.  A more
scientifically natural question is: what is the effect of $T$ on $Y$
that bypasses $M$ when $M$ is held at the value it \emph{would
naturally take} under the reference treatment $T = 0$?

Answering this requires comparing outcomes across two
\emph{counterfactual worlds simultaneously}: the world where $T = 1$
and $M$ is at its natural value under $T = 1$, versus the world where
$T = 1$ but $M$ is held at its natural value under $T = 0$.  Such a
comparison cannot be expressed using the do-operator alone.  It
requires the \emph{nested potential outcomes} notation $Y(t, M(t'))$,
which denotes the outcome that would be observed if $T$ were set to $t$
and $M$ were simultaneously set to the value it would naturally take
if $T$ were $t'$.  These are called \emph{cross-world counterfactuals}
because $t$ and $t'$ may differ, so they refer to two different
hypothetical states of the world simultaneously.

\subsection{Definitions}

\begin{defn}{Natural Direct and Indirect Effects \citep{pearl2001direct}}
\label{w8:defn:nde-nie}
For binary $T \in \{0,1\}$:
\begin{align}
  \mathrm{NDE}
    &= \E\!\left[Y(1, M(0)) - Y(0, M(0))\right],
    \label{w8:eq:nde} \\[4pt]
  \mathrm{NIE}
    &= \E\!\left[Y(1, M(1)) - Y(1, M(0))\right].
    \label{w8:eq:nie}
\end{align}
The \textbf{natural direct effect} (NDE) is the expected change in the
outcome when $T$ shifts from 0 to 1, holding the mediator at the value
it would naturally take under $T = 0$.  The \textbf{natural indirect
effect} (NIE) is the expected change in the outcome due to the shift in
the mediator from $M(0)$ to $M(1)$, holding $T = 1$ fixed.
\end{defn}

\noindent\textbf{Decomposition.}  The total effect decomposes as:
\begin{equation}
  \mathrm{TE} \;=\; \mathrm{NDE} \;+\; \mathrm{NIE}.
  \label{w8:eq:decomp-natural}
\end{equation}

\noindent\textit{Proof.}
\begin{align*}
  \mathrm{NDE} + \mathrm{NIE}
  &= \E[Y(1, M(0)) - Y(0, M(0))]
   + \E[Y(1, M(1)) - Y(1, M(0))] \\
  &= \E[Y(1, M(1)) - Y(0, M(0))] \\
  &= \E[Y(1) - Y(0)]
  = \mathrm{TE}. \qquad\square
\end{align*}

\subsection{Interpreting the Cross-World Nature}

The NDE involves the counterfactual $Y(1, M(0))$: the outcome when $T$
is set to 1 but $M$ is held at the value it would naturally take under
$T = 0$.  This is a \emph{cross-world} quantity because the two
coordinates of the argument refer to different intervention worlds.  It
cannot be observed for any individual, and in general it cannot be
written as a do-expression.

This is not merely a philosophical subtlety.  The cross-world nature
has direct implications for identification: as shown in the next
section, identifying the NDE and NIE requires assumptions that are
strictly stronger than those needed for the CDE or the total effect.

\begin{remark}
In the linear additive SEM~\eqref{w8:eq:sem-m}--\eqref{w8:eq:sem-y}
with no $T \times M$ interaction, the CDE and NDE coincide for all
$m$: $\mathrm{NDE} = \tau'$ and $\mathrm{NIE} = ab$.  The
Baron--Kenny framework therefore implicitly satisfies the cross-world
independence condition needed for this equality.  When there is a
$T \times M$ interaction, $\mathrm{NDE} \neq \mathrm{CDE}$ and the
linear decomposition breaks down: the two concepts correspond to
different scientific questions.
\end{remark}

\begin{remark}[NDE/NIE and the Causal Framework: NPSEM-IE vs.\ FFRCISTG]
\label{w8:rem:ffrcistg}
Whether the cross-world counterfactual $Y(t', M(t))$ is a well-defined
quantity depends on the causal framework adopted, and this has direct
consequences for whether the NDE and NIE are identified estimands.

Under Pearl's NPSEM-IE framework (see Remark~\ref{w1:rem:npsem-ie} in
Chapter~1), independence of the structural errors implies that potential
outcomes and potential mediators across different intervention worlds are
jointly independent.  As a result, $Y(t', M(t))$ is well-defined for
$t \neq t'$, and the NDE and NIE are identifiable under sequential
ignorability.

Under the Finest Fully Randomized Causally Interpretable Structured
Tree Graph (FFRCISTG) framework of \citet{richardson2014ace}, however,
the situation is different.  In FFRCISTG, a causal graph represents a
\emph{single} intervention world at a time; joint counterfactuals of
the form $(Y(t', m), M(t))$ with $t \neq t'$ cannot be realized on a
single graph, because $t$ and $t'$ refer to two distinct intervention
worlds simultaneously.  Consequently, $Y(t', M(t))$ is a cross-world
quantity that FFRCISTG does not define, and within this framework the
NDE and NIE are not identified estimands.

The FFRCISTG resolution is to replace the NDE and NIE with
\emph{separable effects} --- interventionist estimands defined entirely
within a single world by splitting the causal pathway into two
separate interventions \citep{stensrud2022separable,
stensrud2023generalized}.  These separable effects can be identified
without invoking cross-world counterfactuals.  Strikingly, the
identifying formula for separable effects under FFRCISTG coincides
exactly with the mediation formula~\eqref{w8:eq:mediation-formula} for
the NDE and NIE under NPSEM-IE.  This is another instance of the
same-formula, different-interpretation principle encountered in the IV
context (Chapter~7, Remark~\ref{w7:rem:same-formula}): two different
causal frameworks, each with its own structural assumptions, arrive at
the same observed-data expression but assign it different causal
meanings.

For the purposes of this course, we work within the NPSEM-IE framework,
where the NDE and NIE are well-defined under sequential ignorability.
Students who pursue the FFRCISTG literature will find that the
estimation methods developed in Chapter~9 onward remain applicable,
since the identifying formulas are the same.
\end{remark}

% ============================================================
\section{Identification of Natural Effects}
\label{w8:sec:id-natural}
% ============================================================

\subsection{Sequential Ignorability}

Identification of the NDE and NIE requires four conditions, known
collectively as \emph{sequential ignorability}
\citep{imai2010identification}.  Despite the neutral-sounding name,
these conditions are substantially stronger than the ignorability
assumptions used for the total effect or the CDE.  Two of them ---
Assumptions~3 and~4 below --- govern the mediator--outcome
relationship, which is never directly controlled by the investigator.
Sequential ignorability is routinely invoked in applied mediation
analyses, but it is rarely fully plausible in observational settings
and its most critical component cannot be verified from data.  Students
who have internalized the back-door criterion should resist the
temptation to read sequential ignorability as a mild extension of it:
the two assumptions address different sub-problems, and the harder one
has nothing to do with how $T$ was assigned.

\begin{tcolorbox}[defstyle, title={Sequential Ignorability Assumptions}]
\begin{enumerate}
  \item \textbf{No unmeasured treatment--outcome confounding.}
    \[
      \{Y(t', m),\, M(t)\} \;\indep\; T \;\Big|\; \mathbf{X}
      \quad \text{for all } t, t', m.
    \]
    Graphically: $\mathbf{X}$ blocks all back-door paths from $T$ to
    $Y$ and from $T$ to $M$.
  \item \textbf{No unmeasured treatment--mediator confounding.}
    Implied by Assumption~1 above for the $T \to M$ sub-problem.
  \item \textbf{No unmeasured mediator--outcome confounding given $T$.}
    \[
      Y(t', m) \;\indep\; M \;\Big|\; T, \mathbf{X}
      \quad \text{for all } t', m.
    \]
    Graphically: $(T, \mathbf{X})$ blocks all back-door paths from $M$
    to $Y$.
  \item \textbf{No treatment-induced mediator--outcome confounder.}
    There is no variable $L$ such that $T \to L$, $L \to M$ or $L \to Y$,
    and $L$ confounds the $M$--$Y$ relationship.
\end{enumerate}
\end{tcolorbox}

\noindent Assumptions~1--3 are the natural extensions of the Baron--Kenny
conditions to the potential outcomes setting.  Assumption~4 is the
critical new requirement: it rules out a variable $L$ that (i)~is
caused by the treatment and (ii)~confounds the mediator--outcome
relationship.  If such a variable exists, conditioning on $T$ in
Assumption~3 induces a new source of bias that cannot be removed by
further adjustment on observed covariates.

\medskip\noindent\textbf{What randomization of $T$ does and does not
provide.}  Randomizing the treatment $T$ satisfies Assumption~1 (no
unmeasured $T$--$Y$ confounding) and Assumption~2 (no unmeasured
$T$--$M$ confounding) by design.  It does \emph{not} satisfy
Assumption~3 or Assumption~4.  The mediator $M$ is a post-treatment
variable that is never randomized; any unobserved variable $V$ with
$V \to M$ and $V \to Y$ violates Assumption~3 regardless of how $T$
was assigned.  Assumption~4 is even more demanding: it can be violated
by a variable $L$ that is itself \emph{caused} by the treatment, so
randomizing $T$ actually creates the conditions under which
Assumption~4 violations can arise.  The upshot is that a randomized
experiment identifies the total effect and the first-stage $T \to M$
effect cleanly, but it does \emph{not} by itself identify the
second-stage $M \to Y$ effect, and therefore it does not by itself
identify the NDE or NIE.

\begin{warn}{Assumption~4 Is Impossible to Satisfy by Design}
Unlike unmeasured confounders of the $T$--$Y$ relationship (which can
sometimes be eliminated by randomizing $T$), a treatment-induced
mediator--outcome confounder $L$ arises precisely \emph{because} the
treatment changes the system.  Since $T$ causally precedes $M$, any
variable on the path $T \to L \to M$ or $T \to L \to Y$ is a
potential violation of Assumption~4.  This is why identification of
natural effects is considered strictly harder than identification of
the total effect or the CDE.
\end{warn}

\subsection{The Mediation Formula}

Under the four sequential ignorability assumptions, the cross-world
expectation $\E[Y(t, M(t'))]$ is identified from observational data.

\begin{thm}{Mediation Formula \citep{pearl2001direct}}
\label{w8:thm:mediation-formula}
Under sequential ignorability and with discrete $M$ and $\mathbf{X}$:
\begin{equation}
  \E\!\left[Y(t, M(t'))\right]
  = \sum_{m}\sum_{\mathbf{x}}
    \E\!\left[Y \mid T{=}t,\, M{=}m,\, \mathbf{X}{=}\mathbf{x}\right]
    P(M{=}m \mid T{=}t',\, \mathbf{X}{=}\mathbf{x})\,
    P(\mathbf{X}{=}\mathbf{x}).
  \label{w8:eq:mediation-formula}
\end{equation}
\end{thm}

\noindent\textit{Proof sketch.}  By Assumption~3 (no unmeasured $M$--$Y$
confounding given $T$) and Assumption~4 (no treatment-induced
confounder), the outcome $Y(t', m)$ is independent of $M$ given
$(T, \mathbf{X})$:
\begin{align*}
  \E[Y(t, M(t'))]
  &= \sum_{m,\mathbf{x}}
     \E[Y(t, m) \mid M{=}m,\, T{=}t',\, \mathbf{X}{=}\mathbf{x}]
     P(M{=}m \mid T{=}t', \mathbf{X}{=}\mathbf{x})\,
     P(\mathbf{X}{=}\mathbf{x}) \\
  &= \sum_{m,\mathbf{x}}
     \E[Y(t, m) \mid T{=}t,\, M{=}m,\, \mathbf{X}{=}\mathbf{x}]
     P(M{=}m \mid T{=}t', \mathbf{X}{=}\mathbf{x})\,
     P(\mathbf{X}{=}\mathbf{x}) \\
  &= \sum_{m,\mathbf{x}}
     \E[Y \mid T{=}t,\, M{=}m,\, \mathbf{X}{=}\mathbf{x}]
     P(M{=}m \mid T{=}t', \mathbf{X}{=}\mathbf{x})\,
     P(\mathbf{X}{=}\mathbf{x}),
\end{align*}
where the second equality applies Assumption~3, and the third uses
consistency ($Y = Y(t, m)$ when $T = t, M = m$). \hfill$\square$

\medskip\noindent\textbf{Interpretation.}  The mediation formula
``mixes'' the outcome regression under $T = t$ with the mediator
distribution under $T = t'$.  To compute the NDE, set $t = 1$ and
$t' = 0$: take the conditional mean of $Y$ evaluated at treatment 1,
but weight the mediator by its distribution under treatment 0.  This
counterfactual reweighting is what makes the formula non-trivial.

\medskip\noindent\textbf{The NDE and NIE from the formula.}  Applying
\eqref{w8:eq:mediation-formula} with $(t, t') = (1, 0)$ and
$(t, t') = (0, 0)$:
\begin{align*}
  \mathrm{NDE}
  &= \sum_{m,\mathbf{x}}
     \bigl[\E[Y \mid T{=}1, M{=}m, \mathbf{X}{=}\mathbf{x}]
          - \E[Y \mid T{=}0, M{=}m, \mathbf{X}{=}\mathbf{x}]\bigr]
     P(M{=}m \mid T{=}0, \mathbf{X}{=}\mathbf{x})\,
     P(\mathbf{X}{=}\mathbf{x}), \\[4pt]
  \mathrm{NIE}
  &= \sum_{m,\mathbf{x}}
     \E[Y \mid T{=}1, M{=}m, \mathbf{X}{=}\mathbf{x}]
     \bigl[P(M{=}m \mid T{=}1, \mathbf{X}{=}\mathbf{x})
          - P(M{=}m \mid T{=}0, \mathbf{X}{=}\mathbf{x})\bigr]
     P(\mathbf{X}{=}\mathbf{x}).
\end{align*}

\begin{remark}
Three things are worth keeping in view before applying the mediation
formula.
\begin{enumerate}
  \item \textbf{Randomization handles two of the four assumptions, not
    all four.}  Randomizing $T$ secures Assumptions~1 and~2 by
    design.  It leaves Assumptions~3 and~4 entirely open: whether
    there is unmeasured $M$--$Y$ confounding, and whether the
    treatment has induced a variable that confounds $M$ and $Y$, are
    empirical questions that the experimental design does not answer.
  \item \textbf{The formula looks like routine adjustment but is not.}
    The expression $\sum_m \E[Y \mid t, m, \mathbf{x}]\,
    P(M{=}m \mid t', \mathbf{x})$ resembles a standardization formula,
    but it is computing a cross-world expectation $\E[Y(t, M(t'))]$,
    not a do-expression.  Its validity rests on Assumptions~3 and~4,
    which have no graphical counterpart that can be read off the DAG
    from observed variables alone.
  \item \textbf{Sequential ignorability is an untestable assumption
    bundle.}  Unlike treatment ignorability, which can sometimes be
    made credible by design or by a rich covariate set, the
    mediator--outcome ignorability in Assumption~3 and the
    no-treatment-induced-confounder condition in Assumption~4 must be
    defended on subject-matter grounds in every application.  When
    that defense is weak, the CDE --- which requires only the
    back-door criterion --- is the more credible estimand.
\end{enumerate}
\end{remark}

\begin{remark}[Principal Stratification and the Survival Average Causal Effect]
\label{w8:rem:sace}
The NDE and NIE decompose the total effect by causal pathway.  An
alternative approach to effect decomposition --- particularly relevant
when the mediator is an \emph{event} that can preclude observation of
the outcome, such as death --- is \emph{principal stratification}
\citep{frangakis2002principal}.

In this framework, units are classified by their potential mediator
values under each treatment arm: the \emph{principal stratum}
$\{M(1) = m_1, M(0) = m_0\}$ is the subpopulation that would
experience mediator level $m_1$ under $T=1$ and $m_0$ under $T=0$.
This classification is latent --- just as compliance types in the IV
framework (Chapter~7) are unobservable --- but it is well-defined
within the potential outcomes language.

The leading application is \textbf{truncation by death}: when the
mediator $M$ is a survival indicator, units with $M(1) = M(0) = 0$
(those who die under both treatments) cannot be observed for any
outcome $Y$.  The estimand of interest is then the
\textbf{Survival Average Causal Effect} (SACE):
\[
  \mathrm{SACE}
  \;=\;
  \E\!\bigl[Y(1) - Y(0) \;\big|\; M(1) = M(0) = 1\bigr],
\]
the average treatment effect among the stratum of units who would
survive under both treatments.  This estimand is analogous to the LATE
in the IV framework: just as LATE conditions on the complier stratum
$\{T(1) > T(0)\}$, the SACE conditions on the always-survivor stratum
$\{M(1) = M(0) = 1\}$.

Identification of the SACE requires a monotonicity condition --- that
$M(1) \geq M(0)$ for all units (treatment can only help, not hurt,
survival) --- and an assumption about treatment effect homogeneity
across principal strata, analogous to how LATE identification requires
monotonicity on treatment.  The SACE is not identified by the mediation
formula~\eqref{w8:eq:mediation-formula}; it requires separate
identification arguments tailored to the principal stratification
structure.

The principal stratification framework is conceptually distinct from
the NDE/NIE framework.  Where NDE and NIE decompose the \emph{pathway}
of the effect, the SACE asks about the effect within a
\emph{subpopulation} defined by potential mediator values.  In clinical
trials with mortality as an intercurrent event, the SACE is often more
policy-relevant than the NDE, because the NDE involves the cross-world
counterfactual $Y(t', M(t))$ for units who may not survive to have an
outcome.
\end{remark}

\medskip\noindent\textbf{Identification versus estimation.}
Sections~\ref{w8:sec:cde}--\ref{w8:sec:id-natural} have been entirely
about causal estimands and their identification: what quantities are we
trying to learn, and under what structural assumptions are they
expressible as functionals of the observed data distribution?  No
specific model for $P(Y \mid T, M, \mathbf{X})$ or
$P(M \mid T, \mathbf{X})$ was required.  Section~\ref{w8:sec:linear}
turns to a parametric linear model that delivers familiar closed-form
estimators under \emph{additional modeling assumptions} --- linearity,
additivity, and no $T \times M$ interaction --- that go beyond the
identification conditions of §\ref{w8:sec:id-natural}.  Fitting the
regression system does not solve the identification problem; it assumes
it has already been solved.  A student who fits the Baron--Kenny
equations and obtains estimates of $\hat a$, $\hat b$, and $\hat\tau'$
has done estimation; whether those estimates have a causal
interpretation depends entirely on whether the identification
assumptions of §\ref{w8:sec:id-natural} hold in the application.

% ============================================================
\section{The Linear Mediation Model: A Historical Special Case}
\label{w8:sec:linear}
% ============================================================

The framework developed in §\ref{w8:sec:nie}--\ref{w8:sec:id-natural}
--- potential outcomes, cross-world counterfactuals, sequential
ignorability --- is the modern conceptual foundation of mediation
analysis.  Before that framework existed, practitioners used a simpler
regression-based approach that works cleanly under linearity and no
interaction.  That approach, due to \citet{baron1986moderator}, was
enormously influential and remains widely cited.  We study it here for
three reasons: it builds intuition for the two-stage structure of
mediation, its coefficients $a$, $b$, $\tau'$ reappear as special cases
of the NDE and NIE under linearity, and it is the dominant approach in
many applied literatures that students will encounter.

It is not, however, the general framework.  The algebraic
decomposition $\tau = \tau' + ab$ is a consequence of linearity, not a
causal identity.  Outside linear, no-interaction models it fails, and
the product and difference methods it produces are not estimates of the
NDE or NIE.  The goal of this section is to understand what Baron--Kenny
achieves, to be precise about the conditions under which it is valid,
and to be clear about where it breaks down.

\subsection{The Baron--Kenny Three-Equation System}

The regression-based approach of \citet{baron1986moderator} restricts
the prototype graph of Figure~\ref{w8:fig:mediation-dag} to a linear
SEM, conditioning throughout on $\mathbf{X}$:

\begin{align}
  Y  &= \alpha_1 + \tau T + \boldsymbol{\gamma}_1^{\top}\mathbf{X}
        + \varepsilon_1,
  \label{w8:eq:bk1} \\[4pt]
  M  &= \alpha_2 + a T + \boldsymbol{\gamma}_2^{\top}\mathbf{X}
        + \varepsilon_2,
  \label{w8:eq:bk2} \\[4pt]
  Y  &= \alpha_3 + \tau' T + b M
        + \boldsymbol{\gamma}_3^{\top}\mathbf{X} + \varepsilon_3.
  \label{w8:eq:bk3}
\end{align}

The four coefficients of interest are:
\begin{itemize}
  \item $\tau$: the total effect of $T$ on $Y$ (Eq.~\eqref{w8:eq:bk1});
  \item $a$: the first-stage effect of $T$ on $M$ (Eq.~\eqref{w8:eq:bk2});
  \item $\tau'$: the direct effect of $T$ on $Y$ controlling for $M$
    (Eq.~\eqref{w8:eq:bk3});
  \item $b$: the second-stage effect of $M$ on $Y$ controlling for $T$
    (Eq.~\eqref{w8:eq:bk3}).
\end{itemize}

\subsection{The Component Pathways}

Equations \eqref{w8:eq:bk1}--\eqref{w8:eq:bk3} identify three
distinct causal sub-problems, each by the back-door formula.

\medskip\noindent\textbf{First stage} ($T \to M$).  Equation
\eqref{w8:eq:bk2} is the back-door formula for the effect of $T$ on
$M$: conditioning on $\mathbf{X}$ blocks all back-door paths from $T$
to $M$ in $\Gcal$.  In a randomized experiment, $a$ is identified
without conditioning on anything.

\medskip\noindent\textbf{Second stage} ($M \to Y$ given $T$).
Equation \eqref{w8:eq:bk3} implements the back-door formula for the
effect of $M$ on $Y$ given $T$: conditioning on $(T, \mathbf{X})$
blocks the paths $M \leftarrow T \leftarrow U \to Y$ and
$M \leftarrow \mathbf{X} \to Y$.

\begin{warn}{The Second Stage Is Harder Than It Looks}
Even in a randomized experiment where $T$ is randomized, the mediator
$M$ is never randomized.  An unobserved variable $V$ with $V \to M$
and $V \to Y$ creates a back-door path from $M$ to $Y$ that
conditioning on $(T, \mathbf{X})$ cannot block.  This is the central
challenge of mediation analysis: identifying the second-stage effect
$b$ requires a no-unmeasured-confounding assumption for the
mediator--outcome relationship, an assumption that randomisation of
the treatment does not provide.
\end{warn}

Figure~\ref{w8:fig:bk-violation} illustrates the violation.

\begin{figure}[h]
\centering
\begin{tikzpicture}[node distance=1.9cm]
  \node[node]  (T) at (0,0)      {$T$};
  \node[node]  (M) at (2.4,0)    {$M$};
  \node[node]  (Y) at (4.8,0)    {$Y$};
  \node[node]  (X) at (1.2, 1.5) {$\mathbf{X}$};
  \node[unode] (U) at (2.4,-1.5) {$U$};
  \node[unode] (V) at (3.6, 1.5) {$V$};

  \draw[edge]  (T) -- (M);
  \draw[edge]  (M) -- (Y);
  \draw[edge]  (T) to[bend right=20] (Y);
  \draw[edge]  (X) -- (T);
  \draw[edge]  (X) -- (M);
  \draw[dedge] (U) -- (T);
  \draw[dedge] (U) -- (Y);
  \draw[dedge] (V) -- (M);
  \draw[dedge] (V) -- (Y);
  \node[font=\small\itshape\color{darkgrey}, below=0.10cm of V]
    {$M$--$Y$ confounder};
\end{tikzpicture}
\caption{Violation of Assumption~3 (no unmeasured mediator--outcome
confounding).  The unobserved variable $V$ creates a back-door path
$M \leftarrow V \to Y$ that cannot be blocked by conditioning on
$(T, \mathbf{X})$.  Neither randomisation of $T$ nor adjustment for
$\mathbf{X}$ removes this bias.}
\label{w8:fig:bk-violation}
\end{figure}

\subsection{The Product and Difference Formulas}

In the linear SEM~\eqref{w8:eq:bk1}--\eqref{w8:eq:bk3}, the total
effect admits an additive decomposition.

\begin{thm}{Mediation Decomposition in the Linear Model}
\label{w8:prop:decomp}
Under Equations~\eqref{w8:eq:bk1}--\eqref{w8:eq:bk3}:
\begin{equation}
  \tau \;=\; \tau' \;+\; a b.
  \label{w8:eq:linear-decomp}
\end{equation}
The indirect effect and direct effect are therefore identified by:
\begin{align}
  \tau_{\mathrm{ind}} &= ab \quad\text{(\emph{product method})},
  \label{w8:eq:product} \\
  \tau_{\mathrm{dir}} &= \tau - ab = \tau'
    \quad\text{(\emph{difference method})}.
  \label{w8:eq:diff}
\end{align}
\end{thm}

\begin{proof}
Substitute~\eqref{w8:eq:bk2} into~\eqref{w8:eq:bk3}:
\begin{align*}
  Y &= \alpha_3 + \tau' T
       + b(\alpha_2 + aT + \boldsymbol{\gamma}_2^{\top}\mathbf{X}
           + \varepsilon_2)
       + \boldsymbol{\gamma}_3^{\top}\mathbf{X} + \varepsilon_3 \\
    &= (\alpha_3 + b\alpha_2)
       + (\tau' + ab)\,T
       + (\boldsymbol{\gamma}_3 + b\boldsymbol{\gamma}_2)^{\top}\mathbf{X}
       + (b\varepsilon_2 + \varepsilon_3).
\end{align*}
Comparing with~\eqref{w8:eq:bk1} gives $\tau = \tau' + ab$.
\end{proof}

\begin{warn}{The Decomposition $\tau = \tau' + ab$ Is an Algebraic
  Identity, Not a Causal Theorem}
\label{w8:warn:bk-algebraic}
The equality $\tau - \tau' = ab$ is a purely algebraic consequence of
substituting one linear equation into another.  It holds because
linearity makes the indirect effect separable and additive.  It does
\emph{not} hold in general:
\begin{itemize}
  \item In \textbf{nonlinear models} (binary outcomes, count outcomes,
    survival models), the product and difference methods yield
    numerically different estimates.  Neither equals the NIE in general.
  \item When \textbf{$T$ and $M$ interact} in their effect on $Y$,
    the CDE depends on $m$, the NDE and CDE diverge, and the indirect
    effect cannot be summarized by a single number $ab$.
  \item For \textbf{non-continuous mediators}, the product $ab$ has no
    simple causal interpretation outside the linear normal model.
\end{itemize}
Students who learn the Baron--Kenny decomposition first often
overgeneralize it.  The correct generalization is the mediation formula
(Theorem~\ref{w8:thm:mediation-formula}), which reduces to $ab$ and
$\tau'$ only in the linear, no-interaction special case.
\end{warn}

\begin{table}[h]
\centering
\renewcommand{\arraystretch}{1.5}
\begin{tabular}{@{}lll@{}}
\toprule
\textbf{Effect} & \textbf{Formula} & \textbf{Path(s)} \\
\midrule
Total        & $\tau$          & $T \to Y$ and $T \to M \to Y$ combined \\
Direct       & $\tau' = \tau - ab$ & $T \to Y$ only \\
Indirect     & $ab$            & $T \to M \to Y$ only \\
Proportion mediated & $ab/\tau$ & Share of total effect via $M$ \\
\bottomrule
\end{tabular}
\caption{Mediation decomposition in the Baron--Kenny linear model.}
\label{w8:tab:decomp}
\end{table}

\subsection{Inference: The Sobel Test and Bootstrap}

Since $\hat a$ and $\hat b$ are estimated from two separate regressions,
the variance of the product $\hat a \hat b$ is approximated by the
delta method:
\begin{equation}
  \widehat{\mathrm{Var}}(\hat a \hat b)
  \;\approx\;
  \hat b^2\, \widehat{\mathrm{Var}}(\hat a)
  \;+\;
  \hat a^2\, \widehat{\mathrm{Var}}(\hat b).
  \label{w8:eq:sobel}
\end{equation}
The resulting $z$-statistic $z = \hat a \hat b \,/\,
\sqrt{\widehat{\mathrm{Var}}(\hat a \hat b)}$ is the \emph{Sobel test}
\citep{sobel1982asymptotic}.  In practice, bootstrap confidence
intervals for $ab$ are preferred because the distribution of a product
of estimates is skewed in finite samples, especially when either
$\hat a$ or $\hat b$ is small.

\begin{warn}{The ``Significance of Both Paths'' Criterion Is Not a
  Test for Mediation}
\label{w8:warn:sobel}
A common misuse of the Baron--Kenny framework, sometimes called the
``causal steps'' approach, declares mediation to be present when:
(i)~$T \to Y$ is significant in Eq.~\eqref{w8:eq:bk1}; (ii)~$T \to M$
is significant in Eq.~\eqref{w8:eq:bk2}; (iii)~$M \to Y$ is
significant in Eq.~\eqref{w8:eq:bk3}.  This approach has three serious
defects.
\begin{enumerate}
  \item \textbf{Statistical significance $\neq$ mediation.}  A
    significant $\hat a$ and $\hat b$ does not imply that the
    product $ab$ is meaningfully large, nor that it is identified.
    Low power for $\hat a$ or $\hat b$ individually does not imply low
    power for $ab$ (and vice versa).
  \item \textbf{The Sobel test inherits all the assumptions.}  The
    $z$-statistic for $\hat a \hat b$ is only valid under
    Conditions~1--3 and Restrictions~4--5 of the Baron--Kenny
    framework.  In particular it requires no $M$--$Y$ confounding
    (Condition~3) and no $T \times M$ interaction
    (Restriction~5) --- assumptions that are often not checked.
  \item \textbf{Zero total effect does not preclude indirect effects.}
    It is possible for $\tau \approx 0$ while $ab \neq 0$ and
    $\tau' \neq 0$ (direct and indirect effects of opposite sign that
    cancel).  Requiring significance of $\hat\tau$ as a precondition
    will miss these cases entirely.
\end{enumerate}
The modern alternative is to estimate $ab$ (or the NIE directly via
the mediation formula), construct bootstrap or nonparametric confidence
intervals, and interpret the result as a point estimate with
uncertainty --- not as a significance test.
\end{warn}

\subsection{The Baron--Kenny Assumptions: Two Distinct Categories}

The conditions required for the Baron--Kenny decomposition fall into
two fundamentally different categories that should not be conflated.

\begin{tcolorbox}[defstyle,
  title={Baron--Kenny Causal Ignorability Conditions}]
These are causal identification assumptions.  They concern unmeasured
confounding and have graphical interpretations.  Violating them
introduces bias that no amount of additional data can remove given the
observed variables.
\begin{enumerate}
  \item \textbf{No unmeasured $T$--$Y$ confounding.}
    $\varepsilon_1 \indep T \mid \mathbf{X}$.
    Graphically: $\mathbf{X}$ blocks all back-door paths from $T$
    to $Y$.
  \item \textbf{No unmeasured $T$--$M$ confounding.}
    $\varepsilon_2 \indep T \mid \mathbf{X}$.
    Graphically: $\mathbf{X}$ blocks all back-door paths from $T$
    to $M$.
  \item \textbf{No unmeasured $M$--$Y$ confounding given $T$.}
    $\varepsilon_3 \indep M \mid T, \mathbf{X}$.
    Graphically: $(T, \mathbf{X})$ blocks all back-door paths from
    $M$ to $Y$.  This condition is \emph{not} implied by
    randomization of $T$.
\end{enumerate}
\end{tcolorbox}

\begin{tcolorbox}[defstyle,
  title={Baron--Kenny Structural Modeling Restrictions}]
These are parametric modeling assumptions.  They concern the functional
form of the structural equations, not unmeasured confounding.
Violating them does not introduce identification bias in the causal
sense, but it does mean that the product formula $ab$ and the
coefficient $\tau'$ no longer equal the NDE and NIE.
\begin{enumerate}
  \setcounter{enumi}{3}
  \item \textbf{Linearity and additivity.}  The structural equations
    \eqref{w8:eq:bk1}--\eqref{w8:eq:bk3} are correctly specified as
    linear and additive in their arguments.
  \item \textbf{No $T \times M$ interaction.}  The effect of $M$ on
    $Y$ does not depend on the level of $T$, i.e.\ the coefficient
    $b$ in Eq.~\eqref{w8:eq:bk3} is the same for all values of $T$.
\end{enumerate}
Restrictions~4 and~5 can in principle be tested by examining residuals
or adding interaction terms.  Conditions~1--3 cannot be tested from
the observed data alone.
\end{tcolorbox}

\noindent Conditions~1 and~2 hold in Figure~\ref{w8:fig:mediation-dag}
when $\mathbf{X}$ blocks the paths through $U$.  Condition~3 requires
the absence of any variable $V \to M$ and $V \to Y$ not captured by
$(T, \mathbf{X})$; this is never guaranteed by randomization of $T$.
Restrictions~4 and~5 have no graphical counterpart: they must be
justified by domain knowledge or tested by adding a $T \times M$ term
to Eq.~\eqref{w8:eq:bk3} and checking model fit.

\medskip\noindent\textbf{Baron--Kenny as a special case of the modern
framework.}  Under Conditions~1--3 and Restrictions~4--5, the
Baron--Kenny decomposition is a valid estimator of the NDE and NIE
because linearity and no interaction guarantee $\mathrm{NDE} = \tau'$
and $\mathrm{NIE} = ab$.
Outside those conditions, $\tau'$ and $ab$ estimate neither the CDE nor
the natural effects: they estimate regression coefficients in
misspecified structural equations.  The conceptual foundation of
mediation analysis is the potential outcomes framework of
§\ref{w8:sec:nie}--\ref{w8:sec:id-natural}; Baron--Kenny is its
earliest and most restrictive parametric implementation.

\begin{tcolorbox}[defstyle, title={Framework Map}]
\begin{center}
\renewcommand{\arraystretch}{1.4}
\begin{tabular}{@{}p{5.0cm}ll@{}}
\toprule
\textbf{Estimand} & \textbf{Framework needed} & \textbf{Section} \\
\midrule
Total effect $P(y \mid \doop(t))$
  & Do-calculus   & Ch.~5 (back-door) \\
Controlled direct effect (CDE)
  & Do-calculus   & \S\ref{w8:sec:cde} \\
Natural direct effect (NDE)
  & Potential outcomes & \S\ref{w8:sec:nie} \\
Natural indirect effect (NIE)
  & Potential outcomes & \S\ref{w8:sec:nie} \\
Front-door identification
  & Do-calculus   & \S\ref{w8:sec:frontdoor} \\
\bottomrule
\end{tabular}
\end{center}
\end{tcolorbox}

% ============================================================
\section{Front-Door Identification}
\label{w8:sec:frontdoor}
% ============================================================

\subsection{The Front-Door DAG}

Every identification strategy in
§\ref{w8:sec:cde}--\ref{w8:sec:id-natural} assumed that an observed
covariate set $\mathbf{X}$ exists that blocks the back-door paths from
$T$ to $Y$ through $U$.  What if $U$ is wholly unobserved and no such
adjustment set exists?  In the prototype mediation DAG of
Figure~\ref{w8:fig:mediation-dag}, the back-door criterion then fails
for the total effect of $T$ on $Y$, and the methods of earlier sections
are unavailable.

The front-door criterion turns this obstacle into an opportunity: under
two additional restrictions on the prototype mediation graph, the
mediation structure itself provides identification of the total effect
without conditioning on $U$.  The two restrictions are:
\begin{itemize}
  \item remove the direct $T \to Y$ edge, so that $M$ fully mediates
    the effect of $T$ on $Y$; and
  \item require that $U$ has no arrow into $M$, so that the $T \to M$
    sub-effect is unconfounded.
\end{itemize}
The resulting graph is shown in Figure~\ref{w8:fig:frontdoor}.

\begin{figure}[h]
\centering
\begin{tikzpicture}[node distance=1.9cm]
  \node[node]  (T) at (0,0)    {$T$};
  \node[node]  (M) at (2.8,0)  {$M$};
  \node[node]  (Y) at (5.6,0)  {$Y$};
  \node[unode] (U) at (2.8,1.8) {$U$};

  \draw[edge]  (T) -- (M);
  \draw[edge]  (M) -- (Y);
  \draw[dedge] (U) -- (T);
  \draw[dedge] (U) -- (Y);
  \node[font=\small\itshape\color{darkgrey}, below=0.12cm of T] {treatment};
  \node[font=\small\itshape\color{darkgrey}, below=0.12cm of M] {mediator};
  \node[font=\small\itshape\color{darkgrey}, below=0.12cm of Y] {outcome};
  \node[font=\small\itshape\color{darkgrey}, above=0.10cm of U] {unobserved};
\end{tikzpicture}
\caption{The front-door graph.  Compared with the prototype mediation
DAG (Figure~\ref{w8:fig:mediation-dag}), two edges are absent: the
direct $T \to Y$ edge is removed, and $U$ has no arrow into $M$.  These
two omissions are what enable identification: $M$ fully mediates the
effect of $T$ on $Y$, and the $T \to M$ sub-effect is free of
confounding.  The front-door formula identifies $P(y \mid \doop(t))$
without observing or conditioning on $U$.}
\label{w8:fig:frontdoor}
\end{figure}

\begin{tcolorbox}[defstyle, title={Two Different Questions About a Mediator-Like Variable}]
\label{w8:box:fd-vs-mediation}
The front-door graph looks like the mediation graph with two edges
removed, and $M$ sits between $T$ and $Y$ in both.  The resemblance is
structural, but the questions being asked are fundamentally different.

\medskip
\begin{center}
\renewcommand{\arraystretch}{1.5}
\begin{tabular}{p{2.2cm} p{5.0cm} p{5.0cm}}
\toprule
  & \textbf{Ordinary mediation}
  & \textbf{Front-door identification} \\
\midrule
\textbf{Question}
  & How much of the total effect of $T$ on $Y$ flows through $M$?
  & Can $M$ be used to identify the total effect of $T$ on $Y$ despite
    unmeasured $T$--$Y$ confounding? \\
\textbf{Role of $M$}
  & Pathway: carries part of the causal effect
  & Instrument-like relay: routes around confounding \\
\textbf{Target estimand}
  & NDE, NIE (decomposition of the total effect)
  & $P(y \mid \doop(t))$ (the total effect itself) \\
\textbf{What $M$ must satisfy}
  & Sequential ignorability for $M$--$Y$ sub-problem
  & Conditions~1--3 below (graphical, no potential outcomes required) \\
\textbf{Direct $T \to Y$ edge}
  & Present
  & Absent (required) \\
\bottomrule
\end{tabular}
\end{center}

\medskip\noindent The front-door formula does not decompose the total
effect into direct and indirect components.  It identifies the total
effect as a whole, using $M$ as a relay that is unconfounded on the
$T$-side and can be adjusted on the $Y$-side by marginalizing over $T$.
A researcher who wants to know \emph{how much} of the effect goes
through $M$ needs ordinary mediation analysis.  A researcher who cannot
identify the total effect at all because $T$--$Y$ confounding is
unobserved should ask whether the front-door conditions hold.
\end{tcolorbox}

\subsection{The Three Front-Door Conditions}

\begin{tcolorbox}[defstyle, title={Front-Door Conditions}]
A set $\{M\}$ satisfies the \textbf{front-door criterion} for the
effect of $T$ on $Y$ if:
\begin{enumerate}
  \item \textbf{Full mediation.}  All directed paths from $T$ to $Y$
    pass through $M$ (no direct $T \to Y$ edge).
  \item \textbf{No unblocked back-door path from $T$ to $M$.}
    All back-door paths from $T$ to $M$ are blocked (there are no
    unobserved variables that affect both $T$ and $M$).
  \item \textbf{No unblocked back-door path from $M$ to $Y$ given $T$.}
    All back-door paths from $M$ to $Y$ are blocked by conditioning on
    $T$ (there are no unobserved variables that affect both $M$ and $Y$
    without passing through $T$).
\end{enumerate}
\end{tcolorbox}

In Figure~\ref{w8:fig:frontdoor}: Condition~1 holds because there is no
$T \to Y$ edge.  Condition~2 holds because the only back-door paths
from $T$ to $M$ would require a path $T \leftarrow U \to \cdots \to M$,
but $U$ has no arrow into $M$.  Condition~3 holds because the only
back-door path from $M$ to $Y$ is $M \leftarrow T \leftarrow U \to Y$,
which is blocked by conditioning on $T$.

\subsection{Derivation of the Front-Door Formula}

\begin{thm}{Front-Door Formula \citep{pearl1995causal}}
\label{w8:thm:frontdoor}
If $M$ satisfies the front-door criterion for the effect of $T$ on $Y$,
and $P(t) > 0$ for all $t$, then:
\begin{equation}
  P\!\left(y \mid \doop(T{=}t)\right)
  = \sum_m P(m \mid t)
    \sum_{t'} P(y \mid m, t')\, P(t').
  \label{w8:eq:frontdoor}
\end{equation}
\end{thm}

\begin{proof}
We apply the do-calculus in two steps, each exploiting one of the
identification conditions.

\medskip\noindent\textbf{Step 1: Identify the effect of $T$ on $M$.}
By Condition~2, there are no unblocked back-door paths from $T$ to $M$,
so in $\Gcal_{\overline{T}}$ the path $T \leftarrow U \to Y$ is severed
and $T$ has no parents.  By Rule~2 of the do-calculus (or simply
because $T$ satisfies the back-door criterion for $M$ with the empty
set):
\[
  P\!\left(m \mid \doop(T{=}t)\right) = P(m \mid t).
\]

\medskip\noindent\textbf{Step 2: Identify the effect of $M$ on $Y$.}
By Condition~3, conditioning on $T$ blocks all back-door paths from $M$
to $Y$.  Therefore, in $\Gcal_{\overline{M}}$:
\[
  P\!\left(y \mid \doop(M{=}m)\right)
  = \sum_{t'} P(y \mid m, t')\, P(t').
\]

\medskip\noindent\textbf{Combining.}  By Condition~1 (full mediation),
the total effect of $T$ factors through $M$:
\[
  P\!\left(y \mid \doop(T{=}t)\right)
  = \sum_m P\!\left(m \mid \doop(T{=}t)\right)
    P\!\left(y \mid \doop(M{=}m)\right)
  = \sum_m P(m \mid t)
    \sum_{t'} P(y \mid m, t')\, P(t'). \quad\square
\]
\end{proof}

\begin{remark}
The two-step structure of the proof is the do-calculus analogue of the
two-stage reasoning in the mediation formula
(Theorem~\ref{w8:thm:mediation-formula}): first identify the $T \to M$
effect, then the $M \to Y$ effect, and compose.  The mediation formula
accomplishes this via sequential ignorability --- conditional independence
assumptions on potential outcomes.  The front-door formula accomplishes
it via graph structure alone: Condition~2 makes Step~1 unconfounded by
construction, and Condition~3 makes Step~2 unconfounded by marginalizing
over $T$.  No covariate set $\mathbf{X}$ is needed, because the
structural restrictions on the graph substitute for the ignorability
assumptions.
\end{remark}

\subsection{Intuition: Routing Around Confounding}

The front-door formula achieves identification in two steps that each
use only unconfounded variation.

\begin{enumerate}
  \item \textbf{$T$ to $M$:}  There is no confounding on the $T \to M$
    edge ($U$ does not affect $M$), so $P(m \mid t)$ is the causal
    effect of $T$ on $M$.
  \item \textbf{$M$ to $Y$:}  There is back-door confounding on $M \to Y$
    through $T \leftarrow U \to Y$, but it is blocked by summing over
    $T$: $\sum_{t'} P(y \mid m, t')\, P(t')$ marginalises out $T$,
    removing the confounding path.
\end{enumerate}

The key insight is that $U$ confounds $T$ and $Y$ but \emph{not} the
$T \to M$ edge.  The front-door formula exploits this asymmetry to
identify the total causal effect without ever observing or conditioning
on $U$.

\subsection{Front-Door vs.\ Prototype Mediation}

The contrast with Figure~\ref{w8:fig:mediation-dag} is instructive.

\medskip
\begin{center}
\renewcommand{\arraystretch}{1.4}
\begin{tabular}{p{3.5cm} p{4.5cm} p{4.5cm}}
\toprule
\textbf{Feature}
  & \textbf{Prototype mediation DAG}
  & \textbf{Front-door DAG} \\
\midrule
Direct $T \to Y$ edge
  & Present
  & Absent \\
$U$ confounds $T$--$Y$
  & Yes
  & Yes \\
Goal
  & Decompose total effect into direct/indirect
  & Identify total effect despite $T$--$Y$ confounding \\
Requires Assumption~3 ($M$--$Y$ no confounding given $T$)
  & Yes
  & Automatically satisfied \\
Identified by front-door formula
  & No (direct $T \to Y$ breaks Condition~1)
  & Yes \\
\bottomrule
\end{tabular}
\end{center}

% ============================================================
\section{Mediation vs.\ Instrumental Variables}
\label{w8:sec:mediation-vs-iv}
% ============================================================

Mediation analysis and instrumental variables are both frameworks that
reason about intermediate variables on the causal path from $T$ to $Y$.
Yet they differ fundamentally in their purpose, the position of the
intermediate variable, and the assumptions they require.

\subsection{Structural Comparison}

\begin{center}
\renewcommand{\arraystretch}{1.5}
\begin{tabular}{p{3.0cm} p{4.5cm} p{4.5cm}}
\toprule
\textbf{Feature}
  & \textbf{Instrumental Variables}
  & \textbf{Mediation Analysis} \\
\midrule
Position of intermediate variable
  & Pre-treatment ($Z$ precedes $T$)
  & Post-treatment ($M$ follows $T$) \\
Causal role
  & Exogenous source of variation in $T$
  & Pathway through which $T$ affects $Y$ \\
Primary goal
  & Identification of $T \to Y$ effect under confounding
  & Decomposition of $T \to Y$ effect into pathways \\
Key assumption
  & Exclusion: $Z$ affects $Y$ only through $T$
  & Sequential ignorability: no unmeasured $M$--$Y$ confounding given $T$ \\
Estimand
  & LATE (Wald) or ATE (2SLS under homogeneity)
  & NDE, NIE, or CDE \\
Unobserved $T$--$Y$ confounders
  & Permitted (IV routes around them)
  & Must be addressed separately (Assumption~1) \\
Testability
  & Relevance testable; exclusion untestable
  & Sequential ignorability untestable \\
\bottomrule
\end{tabular}
\end{center}

\subsection{The Conceptual Contrast}

The contrast is sharpest in terms of what the intermediate variable
\emph{does} in each framework.

\medskip\noindent In IV, $Z$ is a \emph{handle}: it generates exogenous
variation in $T$ that is free of the back-door confounding path
$T \leftarrow U \rightarrow Y$.  The instrument is valuable precisely
because it is \emph{not} on the causal path from $T$ to $Y$ --- the
exclusion restriction says that $Z$ cannot directly affect $Y$.

\medskip\noindent In mediation analysis, $M$ is a \emph{pathway}: it
transmits part of the causal effect of $T$ to $Y$.  The mediator is
valuable precisely because it \emph{is} on the causal path.  The goal
is to quantify how much of the total effect flows through this pathway.

\begin{tcolorbox}[defstyle, title={The Key Conceptual Distinction}]
\begin{center}
\renewcommand{\arraystretch}{1.4}
\begin{tabular}{p{3cm} p{5cm} p{4.5cm}}
\toprule
& \textbf{IV} & \textbf{Mediation} \\
\midrule
What the intermediate variable does
  & Generates clean variation in $T$ (exogenous source)
  & Carries part of $T$'s causal effect (pathway) \\
Exclusion vs.\ inclusion
  & $Z$ excluded from $Y$'s structural equation
  & $M$ included in $Y$'s structural equation \\
Question answered
  & Does $T$ cause $Y$?
  & How does $T$ cause $Y$? \\
\bottomrule
\end{tabular}
\end{center}
\end{tcolorbox}

\subsection{Can the Same Variable Be Both?}

It is worth asking whether the same variable $M$ could serve as both a
mediator and an instrument.  The answer is: \emph{not in the same
analysis}.  As a mediator, $M$ is on the causal path from $T$ to $Y$
(inclusion required).  As an instrument, $M$ would require the
exclusion restriction --- that it affects $Y$ only through $T$ --- which
contradicts $M$ being a mediator.  The two frameworks make
mutually incompatible structural assumptions about the intermediate
variable.

However, the \emph{front-door identification} formula (§\ref{w8:sec:frontdoor})
is a bridge between the two: it uses the mediator $M$ to achieve
identification of the total effect of $T$ on $Y$ even when $T$ is
confounded --- a goal that IV also pursues.  In this sense, the
front-door formula is IV's closest relative within mediation analysis.

% ============================================================
\section{Summary}
\label{w8:sec:summary}
% ============================================================

\begin{enumerate}
  \item \textbf{Mediation studies mechanisms.}  Mediation analysis
    decomposes the total effect of $T$ on $Y$ into a direct component
    ($T \to Y$) and an indirect component ($T \to M \to Y$).  The
    prototype graph has both pathways; the challenge is that the
    mediator is a post-treatment variable that may be confounded.

  \item \textbf{The total effect is the baseline.}  The TE $= \E[Y(1)
    - Y(0)]$ captures all pathways combined.  It is identified by the
    standard back-door formula; mediation analysis asks for a finer
    decomposition.

  \item \textbf{The CDE uses do-calculus.}  The controlled direct
    effect fixes $M = m$ by joint intervention $\doop(T, M)$ and is
    identified by the back-door formula applied to the pair $(T, M)$.
    No assumptions beyond the back-door criterion are needed, but the
    CDE depends on the fixed level $m$ and does not have a natural
    ``indirect'' complement.

  \item \textbf{Natural effects require potential outcomes.}  The NDE
    and NIE involve cross-world counterfactuals $Y(t, M(t'))$ that
    cannot be expressed with the do-operator alone.  They decompose the
    total effect as TE $=$ NDE $+$ NIE, and are identified by the
    mediation formula under sequential ignorability.  The critical
    Assumption~4 --- no treatment-induced mediator--outcome confounder
    --- is impossible to guarantee by design.

  \item \textbf{The linear model simplifies but restricts.}  The
    Baron--Kenny three-equation system gives the clean formulas
    $\tau_{\mathrm{ind}} = ab$ and $\tau_{\mathrm{dir}} = \tau'$, with
    $\tau = \tau' + ab$.  This decomposition is purely algebraic and
    holds only under linearity and no interaction.  In nonlinear or
    interaction settings, the product and difference methods disagree
    and NDE $\neq$ CDE.

  \item \textbf{Front-door identification uses mediation structure.}
    When $M$ fully mediates $T \to Y$, no $T \to M$ confounding exists,
    and $T$ blocks $M \to Y$ confounding, the front-door formula
    identifies the total effect despite unobserved $T$--$Y$ confounding.
    It achieves this by composing two unconfounded sub-effects:
    $T \to M$ (no confounders) and $M \to Y$ (confounders blocked by
    marginalising over $T$).

  \item \textbf{Mediation and IV are complementary, not equivalent.}
    IV uses a pre-treatment variable to generate exogenous variation
    in $T$; mediation uses a post-treatment variable to decompose the
    causal effect.  The exclusion restriction of IV and the sequential
    ignorability of mediation are mutually incompatible for the same
    intermediate variable.
\end{enumerate}

% ============================================================
\newpage
\section*{Problems}
\addcontentsline{toc}{section}{Problems}
% ============================================================

\begin{enumerate}[label=\textbf{\arabic*.}, leftmargin=*, itemsep=10pt]

  \item \textbf{Identifying the CDE.}
  Consider the DAG: $T \to M$, $T \to Y$, $M \to Y$, $X \to T$,
  $X \to M$, $X \to Y$, with all variables observed.
  \begin{enumerate}[label=(\alph*)]
    \item Write the identification formula for
      $\E[Y \mid \doop(T{=}1), \doop(M{=}m)]$ using the back-door
      criterion for the joint intervention $(T, M)$.
    \item Add an unobserved $U$ with $U \to T$ and $U \to Y$.  Does the
      back-door formula still identify the CDE?  Explain which
      condition, if any, fails.
    \item Instead add $U$ with $U \to M$ and $U \to Y$.  Does the
      back-door formula still identify the CDE?  Explain.
  \end{enumerate}

  \item \textbf{CDE vs.\ total effect.}
  In the prototype graph of Figure~\ref{w8:fig:mediation-dag}, suppose
  $\mathbf{X}$ satisfies the back-door criterion for both the total
  effect and the joint intervention $(T, M)$.
  \begin{enumerate}[label=(\alph*)]
    \item Write expressions for the total effect and the $\mathrm{CDE}(m)$
      using the back-door formula.
    \item Under what graphical condition does $\mathrm{CDE}(m)$ equal the
      total effect for all $m$?  Interpret this condition.
  \end{enumerate}

  \item \textbf{Natural direct and indirect effects.}
  Verify the NDE $+$ NIE $=$ TE decomposition algebraically for
  the linear SEM $M = \alpha T + \eta$,
  $Y = \beta T + \gamma M + \varepsilon$ (no interaction).
  \begin{enumerate}[label=(\alph*)]
    \item Compute $Y(t, M(t'))$ in the linear model.  Show that
      $Y(t, M(t')) = \beta t + \gamma(\alpha t') + \text{noise}$.
    \item Derive $\mathrm{NDE} = \beta$ and $\mathrm{NIE} = \alpha\gamma$
      from the definitions~\eqref{w8:eq:nde}--\eqref{w8:eq:nie}.
    \item Confirm $\mathrm{NDE} + \mathrm{NIE} = \beta + \alpha\gamma
      = \tau_{\mathrm{ATE}}$.
    \item Now suppose a $T \times M$ interaction is added:
      $Y = \beta T + \gamma M + \delta (T \cdot M) + \varepsilon$.
      Compute $\mathrm{CDE}(m)$ and $\mathrm{NDE}$.  Show that
      $\mathrm{NDE} \neq \mathrm{CDE}(m)$ when $\delta \neq 0$.
  \end{enumerate}

  \item \textbf{The Baron--Kenny three-equation system.}
  Consider the prototype graph of Figure~\ref{w8:fig:mediation-dag}
  with $U$ absent and the linear
  SEM~\eqref{w8:eq:bk1}--\eqref{w8:eq:bk3}.
  \begin{enumerate}[label=(\alph*)]
    \item State the four identification assumptions.  For each, give the
      graphical condition in terms of back-door paths.
    \item Derive the equality $\tau = \tau' + ab$ algebraically.
    \item Suppose estimated coefficients are $\hat\tau = 0.50$,
      $\hat a = 0.40$, $\hat b = 0.60$, $\hat\tau' = 0.26$.  Compute
      the indirect effect by both the product and difference methods.
      Do they agree?  Compute the proportion mediated.
    \item An analyst reports $\hat a = 0.40$ with
      $\widehat{\mathrm{SE}}(\hat a) = 0.08$, and $\hat b = 0.60$
      with $\widehat{\mathrm{SE}}(\hat b) = 0.10$.  Compute the Sobel
      standard error for $\hat a \hat b$ using~\eqref{w8:eq:sobel}
      and construct an approximate 95\% confidence interval.
  \end{enumerate}

  \item \textbf{The critical role of Assumption~3.}
  Consider the graph of Figure~\ref{w8:fig:bk-violation}, where an
  unobserved $V$ has $V \to M$ and $V \to Y$.  $T$ is randomized.
  \begin{enumerate}[label=(\alph*)]
    \item Identify all back-door paths from $M$ to $Y$ in this graph.
    \item Can any combination of observed variables $(T, \mathbf{X})$
      block all of these paths?  Explain using $d$-separation.
    \item Suppose an analyst fits Eq.~\eqref{w8:eq:bk3} ignoring $V$
      and obtains $\hat b = 0.80$.  In which direction is $\hat b$
      biased if $V$ has positive effects on both $M$ and $Y$?
    \item State the additional data structure that would be needed to
      identify the second-stage effect nonparametrically.
  \end{enumerate}

  \item \textbf{Front-door identification.}
  Consider the front-door graph of Figure~\ref{w8:fig:frontdoor}.
  \begin{enumerate}[label=(\alph*)]
    \item Verify that the three front-door conditions hold.
    \item Walk through the two-step proof of Theorem~\ref{w8:thm:frontdoor}
      for this graph: identify which do-calculus rule justifies each
      step.
    \item Add a direct edge $T \to Y$ to the graph.  Which front-door
      condition is violated?  Does the formula~\eqref{w8:eq:frontdoor}
      still hold?
    \item Explain why the front-door graph does not require
      Assumption~3 of the Baron--Kenny framework.  Which back-door
      path, present in Figure~\ref{w8:fig:bk-violation}, is absent here?
  \end{enumerate}

  \item \textbf{Mediation vs.\ instrumental variables.}
  A researcher studies the effect of a job training program ($T$) on
  wages ($Y$).  She proposes two intermediate variables:
  (A) motivation ($M_A$), measured after the program starts;
  (B) a lottery that randomly selects applicants for admission ($Z$),
  measured before the program.
  \begin{enumerate}[label=(\alph*)]
    \item For variable~(A): draw the mediation DAG including $M_A$,
      $T$, $Y$, and an unobserved ability variable $U$.  State the
      sequential ignorability assumption needed to identify the NIE
      through $M_A$.  Explain why randomisation of $T$ does not
      automatically satisfy this assumption.
    \item For variable~(B): draw the IV DAG with $Z$, $T$, $Y$, and
      $U$.  State the three IV assumptions.  Explain why the exclusion
      restriction and the ``mediator inclusion'' of mediation analysis
      are mutually incompatible conditions for the same intermediate
      variable.
    \item The researcher argues that $M_A$ (motivation) and $Z$
      (lottery) are both ``intermediate'' variables and that the
      analyzes are interchangeable.  Write a one-paragraph critique of
      this argument, using the distinctions from
      §\ref{w8:sec:mediation-vs-iv}.
    \item Can the front-door formula of §\ref{w8:sec:frontdoor} be
      applied if motivation $M_A$ fully mediates the effect of $T$ on
      $Y$ and $U$ (unobserved ability) does not directly affect $M_A$?
      State the three conditions and assess whether they hold.
  \end{enumerate}

\end{enumerate}

\ifSubfilesClassLoaded{
  \bibliographystyle{plainnat}
  \bibliography{causal_inference}
}{}

\end{document}